\chapter*{Abstract} % no number
\label{abstract}

\addcontentsline{toc}{chapter}{Abstract} % add to index

The increasing adoption of Machine Learning (ML) and Deep Learning (DL) techniques in Network Intrusion Detection Systems (NIDS) has significantly enhanced the capability to detect sophisticated cyber threats that traditional signature-based methods fail to identify. Unlike conventional detection approaches that rely on predefined rules and known attack signatures, ML-based NIDS learn complex patterns from network traffic data, enabling them to identify previously unseen attacks and zero-day exploits. However, this evolution has also exposed such systems to a new class of sophisticated threats known as adversarial machine learning attacks, where intelligent adversaries deliberately manipulate network traffic to evade detection while preserving the malicious functionality of their attacks.

Current adversarial attacks against NIDS predominantly operate in the feature space, manipulating statistical characteristics extracted from network traffic such as packet sizes, inter-arrival times, and flow durations. While these feature-based approaches have demonstrated effectiveness in controlled environments, they suffer from fundamental limitations in practical deployment scenarios. The irreversible and non-differentiable nature of feature extraction processes makes it extremely difficult to translate feature-space perturbations back into valid, replayable network traffic. This creates a significant gap between theoretical adversarial effectiveness and real-world implementability, where protocol constraints, inter-feature dependencies, and network infrastructure requirements must be satisfied.

Furthermore, existing adversarial frameworks focus primarily on manipulating high-level traffic statistics while largely ignoring structural vulnerabilities in packet processing pipelines. Packet fragmentation, a legitimate networking mechanism defined in RFC 791 for handling packets that exceed Maximum Transmission Unit (MTU) constraints, represents a particularly underexplored attack vector. By fragmenting packets into smaller, protocol-compliant segments, attackers can significantly alter traffic characteristics, disrupt feature extraction algorithms, and exploit weaknesses in packet reassembly logic implemented by many NIDS solutions.

This thesis addresses these critical limitations by developing AEGIS (Adversarial Evaluation and Generation framework for Intrusion detection Systems), a comprehensive framework that integrates FRANTIC (FRAgmeNTation adversarIal Component), an advanced packet fragmentation module, with an existing GAN-PSO-based traffic manipulation framework. Unlike purely feature-space approaches, FRANTIC operates at the IPv4 layer to perform controlled, RFC-compliant packet fragmentation that preserves attack semantics while introducing structural perturbations that challenge contemporary NIDS architectures.

The proposed methodology employs a two-phase adversarial generation pipeline that combines the optimization power of existing feature-space manipulation with novel structural transformations. The first phase utilizes a Generative Adversarial Network (GAN) trained on benign traffic distributions combined with Particle Swarm Optimization (PSO) to identify optimal feature perturbations that minimize detection confidence. The second phase applies stochastic IPv4 fragmentation with configurable parameters including fragmentation probability ($\rho$), minimum fragment size (64 bytes), and maximum fragment size (500 bytes), enabling systematic exploration of structural vulnerability space.

FRANTIC implements four core design principles: strict RFC compliance ensuring all fragments conform to IPv4 standards, configurability through adjustable parameters for systematic experimentation, semantic transparency maintaining invisibility to transport and application layers, and stochastic randomization preventing deterministic evasion signatures. The fragmentation algorithm operates probabilistically on eligible packets, randomly sampling fragment sizes within bounded intervals while ensuring proper 8-byte alignment and correct flag handling required by IPv4 specifications.

Comprehensive experimental evaluation was conducted using the CICIDS2019 dataset across three representative ML-based NIDS architectures: Random Forest (RF), Multi-Layer Perceptron (MLP), and LUCID Convolutional Neural Network. The evaluation methodology compared detection performance across three distinct scenarios: baseline unmodified traffic, traffic manipulated using the original GAN-PSO framework, and traffic enhanced with FRANTIC fragmentation.

The experimental results demonstrate dramatic performance degradation across all evaluated models when subjected to fragmentation-enhanced adversarial traffic. Random Forest models experienced the most severe impact, with F1-scores collapsing from 0.9999 (baseline) to 0.0306 (fragmented), representing a 96.9\% performance degradation. False Negative Rates (FNR) increased dramatically from 0.0001 to 0.9837, indicating that nearly 98\% of malicious traffic successfully evaded detection. Multi-Layer Perceptron networks showed similar vulnerability with F1-scores degrading from 0.9987 to 0.1826 (81.7\% degradation) and FNR increasing from 0.0014 to 0.8922. Even the temporal-aware LUCID architecture, designed to capture sequential dependencies in packet flows, experienced severe degradation with F1-scores declining from 0.9757 to 0.1695 (82.6\% degradation) and FNR rising from 0.0466 to 0.9005.

Attack-specific analysis revealed differential vulnerability patterns across malicious traffic categories. Volumetric attacks such as WebDDoS, NTP amplification, and SYN floods achieved near-complete evasion under fragmentation, with FNR values approaching 1.0 across all models. Protocol-constrained attacks like LDAP and PortMap showed greater resilience but still experienced significant degradation. These findings indicate that fragmentation-based evasion exploits fundamental assumptions in feature extraction pipelines rather than model-specific weaknesses.

Comparative analysis between the original manipulator and the fragmentation-enhanced version confirmed that structural perturbations provide substantially greater evasion effectiveness than feature-space manipulation alone. While the original framework achieved moderate success with FNR increases of 10-20\% for most attack types, the addition of fragmentation consistently amplified evasion success across all attack categories and model architectures. This demonstrates that fragmentation operates along a different vulnerability axis than traditional feature manipulation, enabling compound evasion effects that are more difficult to defend against.

The robustness analysis revealed that architectural complexity does not inherently confer adversarial robustness. Traditional machine learning approaches (RF, MLP) and deep learning architectures (LUCID) all exhibited similar vulnerability patterns under fragmentation, suggesting that the effectiveness of this evasion technique stems from exploiting fundamental gaps in traffic representation rather than specific modeling limitations. The results highlight a critical discrepancy between packet-level network reality and the feature-level abstractions used by learning-based detectors.

Beyond adversarial evaluation, this research contributes significant practical applications including dataset augmentation for training more robust NIDS models, penetration testing methodologies for evaluating deployed security systems, and adversarial training frameworks for improving detection resilience. The AEGIS framework enables systematic generation of realistic adversarial traffic that can be directly used for red team exercises, security assessment, and development of next-generation detection capabilities.

The implications of this work extend beyond specific technical contributions to fundamental questions about NIDS architecture and evaluation methodology. The findings demonstrate that high accuracy on clean benchmark datasets provides a false sense of security when systems face realistic adversarial conditions. The research establishes packet fragmentation as a critical vulnerability class that transcends individual model architectures and highlights the necessity for packet reassembly mechanisms integrated before feature extraction in production NIDS deployments.

Future research directions include extending the framework to support IPv6 fragmentation mechanisms, developing real-time adversarial defense systems capable of detecting and mitigating fragmentation-based evasion, investigating hybrid detection architectures that combine signature-based and anomaly-based approaches, and applying explainable AI techniques to understand the precise mechanisms by which fragmentation circumvents detection algorithms.

This thesis provides the cybersecurity community with both offensive capabilities for realistic security testing and defensive insights for developing more robust intrusion detection systems. By bridging the gap between theoretical adversarial machine learning and practical network security, this work contributes to the critical goal of building detection systems capable of operating effectively in hostile, adversarial environments where attackers actively adapt their strategies to circumvent defensive mechanisms. 